% ----------------------------------------------------------
\chapter{Introdução}
% ----------------------------------------------------------

Segundo \citeonline{da2025uso}, a integração de Tecnologias na Atenção Primária à Saúde (APS) tem promovido avanços significativos na gestão e no cuidado de pacientes crônicos. O uso de abordagens cada vez mais inovadoras é impulsionado pela disponibilidade de dados sociodemográficos e clínicos armazenados em sistemas de informação governamentais. A partir deste cenário, o desenvolvimento de ferramentas que proporcionem suporte à decisão clínica surge como uma solução tecnológica para otimizar a triagem e o direcionamento de recursos públicos, especialmente no combate a doenças crônicas não transmissíveis, cujo impacto estrutural e financeiro sobrecarrega o sistema de saúde \cite{bezerra2025morbidade}.

No contexto da ciência de dados aplicada à saúde, algoritmos de aprendizagem de máquina (\textit{Machine Learning}) destacam-se pela capacidade de identificar padrões sutis e modelar relações não lineares em bases de dados complexas, superando as limitações de escores clínicos tradicionais de pontuação fixa \cite{islam2024comprehensive}. A aplicação de técnicas de mineração de dados sobre pesquisas de abrangência nacional permite o treinamento de modelos classificatórios calibrados especificamente para a realidade fenotípica e social da população brasileira. A possível adoção desse tipo de ferramentas no fluxo assistencial não exige apenas alta precisão, mas também transparência algorítmica, mitigando o efeito "caixa-preta"\footnote{Fenômeno em que um modelo de IA oculta seus processos internos de decisão, tornando o caminho entre a entrada de dados e o resultado final incompreensível para os humanos.} característico de modelos de inteligência artificial \cite{salih2025perspective}.

Alinhado a esta ideia, o presente Trabalho de Conclusão de Curso o propõe o desenvolvimento de um modelo classificatório de \textit{Machine Learning} focado na identificação de indivíduos com perfil de risco metabólico associado ao Diabetes  Mellitus tipo 2 (DM2), operando exclusivamente sobre variáveis não invasivas disponíveis no contexto da APS brasileira visando o apoio à decisão. A fundamentação  teórica, a base de dados utilizada e as decisões metodológicas que orientam o desenvolvimento do modelo são apresentadas nas seções subsequentes deste trabalho.